% Options for packages loaded elsewhere
\PassOptionsToPackage{unicode}{hyperref}
\PassOptionsToPackage{hyphens}{url}
\PassOptionsToPackage{dvipsnames,svgnames,x11names}{xcolor}
%
\documentclass[
  10pt,
  a4paper,
]{article}
\usepackage{amsmath,amssymb}
\usepackage{setspace}
\usepackage{iftex}
\ifPDFTeX
  \usepackage[T1]{fontenc}
  \usepackage[utf8]{inputenc}
  \usepackage{textcomp} % provide euro and other symbols
\else % if luatex or xetex
  \usepackage{unicode-math} % this also loads fontspec
  \defaultfontfeatures{Scale=MatchLowercase}
  \defaultfontfeatures[\rmfamily]{Ligatures=TeX,Scale=1}
\fi
\usepackage{lmodern}
\ifPDFTeX\else
  % xetex/luatex font selection
  \setmainfont[]{Liberation Serif}
  \setsansfont[]{Liberation Sans}
  \setmonofont[]{DejaVu Sans Mono}
\fi
% Use upquote if available, for straight quotes in verbatim environments
\IfFileExists{upquote.sty}{\usepackage{upquote}}{}
\IfFileExists{microtype.sty}{% use microtype if available
  \usepackage[]{microtype}
  \UseMicrotypeSet[protrusion]{basicmath} % disable protrusion for tt fonts
}{}
\makeatletter
\@ifundefined{KOMAClassName}{% if non-KOMA class
  \IfFileExists{parskip.sty}{%
    \usepackage{parskip}
  }{% else
    \setlength{\parindent}{0pt}
    \setlength{\parskip}{6pt plus 2pt minus 1pt}}
}{% if KOMA class
  \KOMAoptions{parskip=half}}
\makeatother
\usepackage{xcolor}
\usepackage[margin=22mm]{geometry}
\usepackage{longtable,booktabs,array}
\usepackage{calc} % for calculating minipage widths
% Correct order of tables after \paragraph or \subparagraph
\usepackage{etoolbox}
\makeatletter
\patchcmd\longtable{\par}{\if@noskipsec\mbox{}\fi\par}{}{}
\makeatother
% Allow footnotes in longtable head/foot
\IfFileExists{footnotehyper.sty}{\usepackage{footnotehyper}}{\usepackage{footnote}}
\makesavenoteenv{longtable}
\setlength{\emergencystretch}{3em} % prevent overfull lines
\providecommand{\tightlist}{%
  \setlength{\itemsep}{0pt}\setlength{\parskip}{0pt}}
\setcounter{secnumdepth}{5}
\ifLuaTeX
\usepackage[bidi=basic]{babel}
\else
\usepackage[bidi=default]{babel}
\fi
\babelprovide[main,import]{korean}
\ifPDFTeX
\else
\babelfont{rm}[]{Liberation Serif}
\fi
% get rid of language-specific shorthands (see #6817):
\let\LanguageShortHands\languageshorthands
\def\languageshorthands#1{}
\usepackage{xeCJK}
\setCJKmainfont{Noto Serif CJK KR}
\setCJKsansfont{Noto Sans CJK KR}
\setCJKmonofont{Noto Sans Mono CJK KR}
\xeCJKsetup{CJKspace=true}
\usepackage{xurl}
\usepackage{fvextra}
\usepackage{fancyhdr}
\AtBeginDocument{\DefineVerbatimEnvironment{verbatim}{Verbatim}{breaklines=true,breakanywhere=true,fontsize=\small}}
\setlength{\emergencystretch}{4em}
\setlength{\headheight}{15pt}
\pagestyle{fancy}\fancyhf{}
\fancyhead[L]{\small 중성자별 구조와 비방사 진동}
\fancyfoot[C]{\thepage}
\renewcommand{\headrulewidth}{0.3pt}
\setcounter{tocdepth}{2}
\ifLuaTeX
  \usepackage{selnolig}  % disable illegal ligatures
\fi
\usepackage{bookmark}
\IfFileExists{xurl.sty}{\usepackage{xurl}}{} % add URL line breaks if available
\urlstyle{same}
\hypersetup{
  pdftitle={설명에서 독립적인 판단으로},
  pdflang={ko-KR},
  colorlinks=true,
  linkcolor={Maroon},
  filecolor={Maroon},
  citecolor={Blue},
  urlcolor={Blue},
  pdfcreator={LaTeX via pandoc}}

\title{설명에서 독립적인 판단으로}
\usepackage{etoolbox}
\makeatletter
\providecommand{\subtitle}[1]{% add subtitle to \maketitle
  \apptocmd{\@title}{\par {\large #1 \par}}{}{}
}
\makeatother
\subtitle{지도와 피드백을 위한 안내}
\author{}
\date{}

\begin{document}
\maketitle

{
\hypersetup{linkcolor=}
\setcounter{tocdepth}{3}
\tableofcontents
}
\setstretch{1.12}
\section{설명이 학습으로 이어지는지
확인하기}\label{uxc124uxba85uxc774-uxd559uxc2b5uxc73cuxb85c-uxc774uxc5b4uxc9c0uxb294uxc9c0-uxd655uxc778uxd558uxae30}

좋은 설명을 읽고 고개를 끄덕이는 일과, 새로운 입력에서 같은 원리를
사용하는 일은 다르다. 그렇다고 설명을 줄이고 처음부터 혼자 풀게 하면
되는 것도 아니다. 이 안내서는 두 극단 사이에서 어떤 도움을 언제 주고,
어떤 증거로 다음 단계에 갈지 정한다. 대상은 TOV 구현 경험과 기본 Cowling
공부가 있는 독자다. 실제 학생에게 효과를 측정한 기록은 아직 없다.

\subsection{이 자료의
역할}\label{uxc774-uxc790uxb8ccuxc758-uxc5eduxd560}

물리 본문은 정의·유도·그림·한계를 담고, 큐시트는 그 지식을 실제로 사용할
한 번의 작업을 정한다. 이 안내서는 지도자가 참고할 근거와 피드백 사례다.
학생에게 고등교육학 문헌을 추가로 공부하도록 요구하지 않는다. 학생이 볼
평가 기준은 \texttt{docs/ASSESSMENT.md}에 따로 둔다.

완성 수치 방법의 호출을 허용하는 것은 선택한 학습목표를 좁힌 것이다. 이
과정의 기본 성취는 범용 적응 적분기나 역함수를 새로 작성하는 능력이
아니다. 입력 물리와 경계조건을 정하고, 제공 함수의 반환값을 검사하며,
수렴 자료를 근거로 제한된 결론을 내리는 능력이다. 직접 구현을 선택한
학생에게는 별도 코딩 증거를 받는다.{[}E1{]}

\subsection{문장의 친절함과 생각의 부담을
구분한다}\label{uxbb38uxc7a5uxc758-uxce5cuxc808uxd568uxacfc-uxc0dduxac01uxc758-uxbd80uxb2f4uxc744-uxad6cuxbd84uxd55cuxb2e4}

본문의 질문은 장식이 아니다. 독자가 막힐 지점을 먼저 드러내고 바로
필요한 식이나 작업으로 넘어가기 위해 쓴다. 하나의 긴 문장에
정의·해석·한계를 한꺼번에 넣지 않는다. 식 뒤 첫 문장은 기호를 정의하거나
직접적인 결과 하나를 말하고, 다음 문장에서 물리적 역할과 조건을
분리한다.

설명을 따라갈 수 없다는 응답은 실제 정보다. 이를 곧바로 학생의 노력
부족으로 해석하지 않는다. 다른 한편으로 설명이 유려해서 이해한 느낌이
들었다는 응답도 독립 수행을 대신하지 않는다. 대학 물리의 연구에서
학습감과 실제 평가 성과가 구별되었지만, 그것을 이유로 글을 어렵게
만들거나 읽기 불편을 무시할 수는 없다.{[}E7{]}

\section{설명·시도·피드백의
배치}\label{uxc124uxba85uxc2dcuxb3c4uxd53cuxb4dcuxbc31uxc758-uxbc30uxce58}

\subsection{새 관계를 처음 배울
때}\label{uxc0c8-uxad00uxacc4uxb97c-uxcc98uxc74c-uxbc30uxc6b8-uxb54c}

새 개념에서는 완성 예제를 먼저 충분히 제공한다. 다만 풀이를 읽는
학생에게 어느 보존법칙이나 정의가 다음 줄을 허용하는지 설명하게 한다.
처음부터 모든 줄을 가린 채 계산시키지 않는다. 역학 예제의 자기설명
연구가 보여 준 것은 풀이 행동과 원리의 연결이며, 말의 양 자체가 성취인
것은 아니다.{[}E2{]}

예를 들어 A09에서는 Eulerian·Lagrangian 변화를 본문에서 설명한 뒤,
표면에서 어느 변화가 0이어야 하는지 말하게 한다. 답이 불분명하면 유체
요소를 옮긴 그림부터 다시 본다. 중심과 표면의 고급 연결 코드를 작성할 수
있느냐는 그 다음 질문이며 기본 성취가 아니다.

\subsection{익숙한 관계를 다음 문제에 사용할
때}\label{uxc775uxc219uxd55c-uxad00uxacc4uxb97c-uxb2e4uxc74c-uxbb38uxc81cuxc5d0-uxc0acuxc6a9uxd560-uxb54c}

풀이의 마지막 판단 하나를 학생에게 맡긴다. 처음의 완전한 예제, 한 단계가
비어 있는 예제, 도움 없이 선택할 새 입력을 순서대로 둔다. 원리를
설명하게 하는 질문과 함께 도움을 줄인 실험은 전이를 지원했지만, 도움을
줄였다는 사실만으로 먼 전이가 보장되지는 않았다.{[}E3{]}

A13에서는 endpoint 조건을 만족하는 cubic을 보여 준 다음 구간 내부의
허용조건을 학생이 고르게 한다. A20에서는 제공 함수를 호출하고 구간 폭이
바뀌었을 때 무엇을 확인할지 결정한다. 두 과제의 연결은 같은 식을
반복하는 데 있지 않고, 그 식을 다른 판단에 사용하는 데 있다.

\subsection{시간이 지난 뒤 다시 물을
때}\label{uxc2dcuxac04uxc774-uxc9c0uxb09c-uxb4a4-uxb2e4uxc2dc-uxbb3cuxc744-uxb54c}

다음 만남에 짧게 답하게 한 뒤 자료를 열어 수정한다. 최초 예측은 채점하지
않으며, 지연 질문의 목적도 기억에서 공식 전체를 재현하는 것이 아니다.
같은 정의가 다른 입력에서 여전히 필요한 이유를 확인한다. 산문 회상에서
나타난 지연 효과는 이 선택의 참고이지 복잡한 문제 해결 전반의 보증은
아니다.{[}E4{]}

반대 결과도 고려한다. 과학 선수지식이 제한적인 대학생의 회로 문제에서는
반복 예제 학습이 예제 뒤 유사 문제를 푸는 조건보다 지연 검사에서
좋았다.{[}E8{]} 그러므로 막힌 학생에게 설명 없이 더 오래 풀라고 요구하지
않는다. 첫 힌트는 정의를 가리키고, 둘째 힌트는 작은 사례를 제시한다.
그래도 어려우면 한 단계를 함께 풀어 준 뒤 다른 입력에서 다시 시도한다.

\subsection{피드백이 다음 행동으로 이어질
때}\label{uxd53cuxb4dcuxbc31uxc774-uxb2e4uxc74c-uxd589uxb3d9uxc73cuxb85c-uxc774uxc5b4uxc9c8-uxb54c}

형성평가의 목적은 점수를 한 번 붙이는 데 있지 않다. 기준을 공개하고
학생이 자신의 결과를 비교·수정할 수 있게 한다.{[}E5{]} 수렴표를 읽은 뒤
다음 실험을 선택하는 활동은 물리 실험에서 연구된 비교와 행동의 반복을 이
계산과정에 맞게 옮긴 것이다. 원 연구는 다른 연도의 집단을 비교했고 측정
불확실성을 다뤘다. 여기의 결정론적 수치 여유를 통계적 신뢰구간으로
바꾸지 않는다.{[}E6{]}

지도자는 불일치마다 tolerance 감소를 처방하지 않는다. 단위·계수·배경
표현·고유값 탐색·모드 식별 중 어떤 증거가 부족한지 먼저 학생에게 묻는다.
그 다음 실제로 가능한 작은 실험 하나를 고른다. 이미 통과한 무관한 계산을
모두 다시 하게 하는 것은 이 활동의 목적이 아니다.

\section{피드백의 구체적인
예}\label{uxd53cuxb4dcuxbc31uxc758-uxad6cuxccb4uxc801uxc778-uxc608}

다음은 가상 응답에 대한 지도 예시다. 실제 학생의 응답이나 측정 자료가
아니다.

\subsection{A01: 왕복 시험}\label{a01-uxc655uxbcf5-uxc2dcuxd5d8}

가상 응답: ``원래 값으로 돌아왔으므로 단위가 맞다.''

바로 정답을 다시 말하기보다 잘못된 순변환 계수의 역수를 역변환에 썼을
때를 묻는다. 학생이 왕복 성공과 독립 단위 정의를 구별하면 다른 단위의 한
값을 직접 변환하게 한다. 증거는 독립 기준값과 그 지수의 설명이다.

\subsection{A09: 한 trial의 해}\label{a09-uxd55c-trialuxc758-uxd574}

가상 응답: ``끝까지 적분했고 잔차도 작으니 모드다.''

먼저 진폭을 작게 바꾸면 정규화하지 않은 잔차가 어떻게 되는지 묻는다. 그
다음 최종 주파수에서의 재적분과 고유함수 식별을 구별한다. 식을
이해하지만 함수가 실패했다면 작은 한 점의 우변으로 내려가 실제 전달값을
확인한다. 두 상황에 같은 처방을 하지 않는다.

\subsection{A13·A20: 부드러운
곡선}\label{a13a20-uxbd80uxb4dcuxb7ecuxc6b4-uxace1uxc120}

가상 응답: ``끝점의 압력과 기울기가 맞으니 물리적으로 허용된다.''

구간 내부의 둘째 미분을 아직 확인하지 않았다는 점을 찾게 한다. 제공된
거절 예제에서 어느 검사에 걸렸는지 읽고, 미분의 knot와 물질적 밀도
jump를 구별하게 한다. 허용조건을 통과시키려고 표를 임의로 평활화하는
답은 원래 모형을 바꾼 선택임을 밝히게 한다.

\subsection{A19·A21: 작은 주파수
차이}\label{a19a21-uxc791uxc740-uxc8fcuxd30cuxc218-uxcc28uxc774}

가상 응답: ``두 별의 차이가 일정하므로 각 주파수도 정확하다.''

두 절대 주파수가 같이 이동하는 경우를 설명하게 한다. 이후 실제
수렴표에서 각각의 변화와 차이의 변화를 따로 찾게 한다. 다음 실험은
학생이 선택하고, 그 실험이 어떤 오류를 배제하는지 말하게 한다. 예상한
부호를 얻었는지는 성취 기준이 아니다.

\subsection{P4: 동료의 실행}\label{p4-uxb3d9uxb8ccuxc758-uxc2e4uxd589}

가상 응답: ``내 컴퓨터에서는 된다.''

그 사실을 부정할 필요는 없다. 동료가 가진 입력과 실행 위치에서 무엇이
달라졌는지 찾아 안내문 한 문장을 고친다. 새 안내문으로 작은 입력을 다시
실행한 기록을 받는다. 구두 설명을 반복하는 대신 빠진 조건을 문서에
남긴다.

\section{목표와 제출 증거의
대응}\label{uxbaa9uxd45cuxc640-uxc81cuxcd9c-uxc99duxac70uxc758-uxb300uxc751}

아래 표는 자동시험을 대체하는 추가 시험 목록이 아니다. 각 과제에서 코드
출력 옆에 어떤 짧은 설명을 받아야 할지 정한다. 모든 항목에 긴 보고서를
요구하지 않는다.

\begin{longtable}[]{@{}
  >{\raggedright\arraybackslash}p{(\columnwidth - 4\tabcolsep) * \real{0.3333}}
  >{\raggedright\arraybackslash}p{(\columnwidth - 4\tabcolsep) * \real{0.3333}}
  >{\raggedright\arraybackslash}p{(\columnwidth - 4\tabcolsep) * \real{0.3333}}@{}}
\toprule\noalign{}
\begin{minipage}[b]{\linewidth}\raggedright
과제
\end{minipage} & \begin{minipage}[b]{\linewidth}\raggedright
학생이 할 판단
\end{minipage} & \begin{minipage}[b]{\linewidth}\raggedright
최소 설명·수행 증거
\end{minipage} \\
\midrule\noalign{}
\endhead
\bottomrule\noalign{}
\endlastfoot
A00 & 프로그램이 끝났다면, 우리는 무엇을 확인한 것일까? & 도구 판본과 첫
명령 출력; 빌드 성공이 아직 증명하지 않은 것 두 가지 \\
A01 & 왕복 변환이 맞으면 단위도 맞은 것일까? & 환산 유도와 차원표; 독립
기준값과 왕복 결과 \\
A02 & 작은 잔차는 근에 가까워졌다는 뜻일까? & bracket 폭과 반복 횟수 표;
성공 사례와 실패 상태 한 개 \\
A03 & 적분값 하나가 맞았는데도 왜 다른 예제를 푸는가? & 단계표·수렴 표
또는 그림; 비자율 방정식 검사 \\
A04 & 거절한 단계는 어디로 돌아가는가? & 거절 전후 trace와 accepted
profile; 크기가 다른 두 성분의 norm 검사 \\
A05 & 압력곡선 하나가 총에너지까지 결정하는가? & 유도와 domain 조건; 세
차분 간격의 μ 비교 \\
A06 & TOV를 이미 풀었는데 왜 중심과 표면을 다시 읽는가? & 중심급수 유도;
M,R 및 μexp(ν) profile \\
A07 & 같은 중심압력으로 비교하면 같은 질량의 별인가? & Pc--M scan; 최종
질량 잔차와 branch 식별 \\
A08 & 부호가 바뀌는 곳마다 고유값이 있는가? & 잔차 scan 및 pole 반례; 두
고유함수와 k\_n L 비교 \\
A09 & 한 번 끝까지 적분한 해는 이미 진동 모드인가? & 차원표와 point RHS;
한 trial의 W,V profile \\
A10 & 두 번째로 찾은 근은 반드시 첫 압력모드인가? & root\_status와
identity\_status; 같은 좌표/정규화의 f,p1 고유함수 \\
A11 & 고밀도 부분을 고정했다는 말은 어디서 확인하는가? & C\_L 유도와
units; 좌우 상태 표 \\
A12 & 압력곡선의 교점 하나로 상전이가 결정되는가? & analytic crossing과
수치 잔차; Δε−μtΔn 검사 \\
A13 & 끝점이 매끄럽게 이어지면 물리적으로도 허용되는가? & cubic 계수와
endpoint residual; admitted/rejected 각각 한 사례 \\
A14 & 밀도가 뛰면 변위의 두 성분도 함께 뛰어야 하는가? & 전달식 유도; W
및 ΔP 좌우값 \\
A15 & 포화조건 세 개는 함수형까지 정해 주는가? & 계수 유도와 실제
calibration residual; 서로 다른 regulator 비교 \\
A16 & 밀도를 바꿀 때 조성도 다시 풀어야 하는가? & 한 밀도의
x\_p,α,μ\_e,잔차; 전하중성 검사 \\
A17 & 가장 낮은 후보는 곧 최소점인가? & 네 에너지 항 분해; coarse/fine
최소 후보 \\
A18 & 표를 읽었다면 별에 넣을 EOS도 준비된 것일까? & raw 열 정의와 SI
변환; admitted/rejected interval 설명 \\
A20 & 같은 표에서 읽은 값과 미분이 왜 서로 어긋날까? & 한 구간의
P,n,epsilon,chi 유도; local t 역함수와 절대 mu subtraction 비교 \\
A19 & 차이가 일정하면 두 계산도 수렴한 것일까? & 핵심 자동시험; 세
수준의 차이·floor·판정 표 \\
A21 & 같은 문제의 해라는 주장을 무엇으로 방어하는가? & 같은 물리 BC와 두
방향 해의 연결; 두 별의 절대 주파수와 paired difference 수렴표 \\
P0 & 옮긴 것은 실행파일인가, 자신이 작성한 계산인가? & 선택한 경로와
SOURCE\_FILES.json 또는 CALCULATION.json; 전후 소스 비교 및 직접
작성·제공 함수의 구분 \\
P1 & 공용 함수를 가져오면 외부 의존성도 사라지는가? & learner export의
실제 의존 경로 검사; 검사한 폴더와 실행 경로의 명시 \\
P2 & 지난 실행의 흔적이 없어도 다시 만들 수 있는가? & 실제 소스 목록과
binary·외부 경로 검사; 검사가 배제한 의존성과 남은 한계 \\
P3 & 빌드가 끝나면 물리 검증도 끝나는가? & 소스부터 빌드한 명령과
environment 실행; 물리 검증과 구분한 성공 범위 \\
P4 & 동료는 설명 없이도 같은 질문을 계산할 수 있는가? & 동료가 사용한
입력·명령·결과; 질문을 반영해 수정한 안내문 한 문장 \\
P5 & 음향시간을 추가하면서 원래 별을 바꿔야 하는가? & cumulative τ CSV와
metadata; 두 독립 적분 기준의 시험 결과 \\
\end{longtable}

\section{실제 수업에서의
확인}\label{uxc2e4uxc81c-uxc218uxc5c5uxc5d0uxc11cuxc758-uxd655uxc778}

\subsection{첫 적용은 형성평가로
한다}\label{uxccab-uxc801uxc6a9uxc740-uxd615uxc131uxd3c9uxac00uxb85c-uxd55cuxb2e4}

자료가 좋아졌다는 제작자의 판단과 실제로 학생이 배웠다는 증거는 다르다.
첫 적용에서는 다음 세 지점을 우선 관찰한다. A01의 단위·정의 구분, A09의
경계조건·한 trial 해석, A21의 절대 주파수·차이·한계 설명이다. 서로 다른
인지 부담을 가진 지점이므로 하나의 완료율로 합치지 않는다.

각 지점에서 최초 답, 어떤 도움을 받았는지, 수정한 답, 새로운 입력에서의
수행, 다음 만남의 짧은 응답을 기록한다. 소요시간에는 설치
대기·실행시간과 실제 개념 작업시간을 구분한다. 도움의 수준은 0=없음,
1=정의 위치, 2=작은 사례, 3=한 단계 공동 풀이, 4=전체 풀이 확인으로 적을
수 있다. 이 구분은 수업 기록을 위한 선택이며 표준화된 심리 척도가
아니다.

\subsection{관찰 기록 한
장}\label{uxad00uxcc30-uxae30uxb85d-uxd55c-uxc7a5}

\begin{longtable}[]{@{}ll@{}}
\toprule\noalign{}
항목 & 기록 \\
\midrule\noalign{}
\endhead
\bottomrule\noalign{}
\endlastfoot
과제·날짜·선택 경로 & \\
처음에 설명한 정의·예측 & \\
막힌 지점과 실제 입력 & \\
제공한 도움의 내용·수준 & \\
수정한 이유 & \\
새 입력에서 유지된 판단 & \\
다음 만남의 응답·필요한 도움 & \\
개념 작업시간·실행 대기시간 & \\
다음에 고칠 문장 또는 과제 하나 & \\
\end{longtable}

별도의 긴 설문보다 이 기록과 실제 응답을 먼저 본다. 읽기 부담을 물을
때에는 어떤 문장에서 무엇을 알 수 없었는지 구체적으로 묻는다. 학생이
어렵다고 했다는 사실을 교육적 효과의 증거로 바꾸지 않는다.

\subsection{이 자료를 다시 고칠
조건}\label{uxc774-uxc790uxb8ccuxb97c-uxb2e4uxc2dc-uxace0uxce60-uxc870uxac74}

여러 번 같은 용어의 뜻을 묻는다면 그 정의의 위치를 옮긴다. 공식을
재현하지만 다른 입력에 적용하지 못하면 설명을 더 외우게 하지 않고 대조
사례를 바꾼다. 도움을 줄인 뒤에도 원리를 선택하면 같은 유형의 기초
연습을 줄인다. 특정 질문이 항상 전제를 설명하는 긴 대화로 바뀌면, 그
전제를 본문이나 큐시트의 앞부분에 넣는다.

학생 한 명의 전후 변화는 사용성·학습 병목을 찾는 데 쓸 수 있지만 자료의
인과효과를 일반화하기에는 부족하다. 문제 친숙도, 지도 시간, 연습량과
입력 난이도가 함께 달라질 수 있다. 더 체계적인 평가를 하려면 동일한
목표를 평가하는 다른 문제와 도움 기록을 사용하고, 가능한 경우 기준을
모르는 두 평가자가 일부 응답을 비교한다. 여기서는 그런 실증을 이미
수행했다고 주장하지 않는다.

\section{문헌과 적용의
경계}\label{uxbb38uxd5ccuxacfc-uxc801uxc6a9uxc758-uxacbduxacc4}

\subsection{{[}E1{]} Biggs, John (1996)}\label{e1-biggs-john-1996}

Enhancing teaching through constructive alignment. \emph{Higher
Education} 32(3), 347--364. DOI:
\url{https://doi.org/10.1007/BF00138871}.

\textbf{확인한 근거.} 교육목표를 실제 수행으로 정의하고 학습 활동과
평가를 그 수행에 맞추는 개념적 틀이다. 특정 중성자별 강의의 학습 효과를
측정한 연구는 아니다.

\textbf{이 강의에서의 선택.} 각 과제에서 무엇을 설명해야 하는지와 실제
제출 증거를 대응시킨다. 자동시험은 그 증거의 일부이며 물리적 설명을
대신하지 않는다.

\subsection{{[}E2{]} Chi, Michelene T. H.
(1989)}\label{e2-chi-michelene-t.-h.-1989}

Self-Explanations: How Students Study and Use Examples in Learning to
Solve Problems. \emph{Cognitive Science} 13(2), 145--182. DOI:
\url{https://doi.org/10.1207/s15516709cog1302_1}.

\textbf{확인한 근거.} 역학 예제를 공부하며 말한 설명을 분석했다. 이해가
좋은 학생은 풀이의 행동을 원리와 적용조건에 더 자주 연결했다. 이
관찰만으로 자기설명 문구의 독립적 인과효과를 정한 것은 아니다.

\textbf{이 강의에서의 선택.} 식의 다음 줄을 베끼게 하기보다 왜 이 정의나
보존법칙을 쓰는지 묻는다. 설명이 짧아도 사용한 조건을 정확히 밝히면
충분하다.

\subsection{{[}E3{]} Atkinson, Robert K.
(2003)}\label{e3-atkinson-robert-k.-2003}

Transitioning From Studying Examples to Solving Problems: Effects of
Self-Explanation Prompts and Fading Worked-Out Steps. \emph{Journal of
Educational Psychology} 95(4), 774--783. DOI:
\url{https://doi.org/10.1037/0022-0663.95.4.774}.

\textbf{확인한 근거.} 두 실험에서 풀이 단계의 점진적 제거와 원리를
설명하게 하는 질문을 결합했을 때 가까운 전이와 먼 전이의 수행이
개선되었다. 풀이를 줄이는 것만으로 먼 전이가 언제나 확보되는 것은
아니었다.

\textbf{이 강의에서의 선택.} 처음에는 완전한 풀이를 주고, 다음에는
마지막 판단 하나를 맡긴다. 이후 새 입력에서 필요한 검사를 스스로 고르게
한다. 구체적 물리 관계의 설명을 함께 요구한다.

\subsection{{[}E4{]} Roediger, Henry L., III
(2006)}\label{e4-roediger-henry-l.-iii-2006}

Test-Enhanced Learning: Taking Memory Tests Improves Long-Term
Retention. \emph{Psychological Science} 17(3), 249--255. DOI:
\url{https://doi.org/10.1111/j.1467-9280.2006.01693.x}.

\textbf{확인한 근거.} 산문을 학습한 학생의 회상 실험에서 반복 읽기는
짧은 지연에, 이전 회상 시험은 더 긴 지연에 유리했다. 이 결과를 복잡한
유도의 전이나 코딩 능력 전체에 그대로 적용할 수는 없다.

\textbf{이 강의에서의 선택.} 다음 만남에 짧게 개념을 복원한 뒤 자료를
열어 수정한다. 유도 전체를 기억에서 재현하도록 강요하지 않고 판단에
필요한 정의와 조건을 다시 사용한다.

\subsection{{[}E5{]} Nicol, David J.
(2006)}\label{e5-nicol-david-j.-2006}

Formative assessment and self-regulated learning: a model and seven
principles of good feedback practice. \emph{Studies in Higher Education}
31(2), 199--218. DOI: \url{https://doi.org/10.1080/03075070600572090}.

\textbf{확인한 근거.} 형성평가와 피드백 연구를 자기조절 학습의 관점에서
정리한 모형이다. 학생을 피드백의 수동적 수신자가 아닌 자기평가와 수정의
주체로 본다. 단일 무작위 실험의 효과 추정치는 아니다.

\textbf{이 강의에서의 선택.} 채점 기준을 공개하고 피드백 뒤의 재시도를
마련한다. 점수만 적지 않고 다음에 무엇 하나를 바꾸어 확인할지 학생과
정한다.

\subsection{{[}E6{]} Holmes, N. G. (2015)}\label{e6-holmes-n.-g.-2015}

Teaching critical thinking. \emph{Proceedings of the National Academy of
Sciences} 112(36), 11199--11204. DOI:
\url{https://doi.org/10.1073/pnas.1505329112}.

\textbf{확인한 근거.} 대학 기초 물리 실험에서 정량적 비교에 따라 행동을
결정하는 활동을 반복하고 지시를 줄였다. 이후 자발적 비교·방법 수정의
차이가 관찰되었다. 다른 연도의 집단을 비교했으므로 무작위 배정 연구로
부르지 않는다.

\textbf{이 강의에서의 선택.} 수렴표를 본 뒤 다음 실험을 선택하고 실제로
수행하게 한다. 측정 불확실성을 다룬 원 연구의 통계량을 결정론적 수치
오차에 옮기지 않는다.

\subsection{{[}E7{]} Deslauriers, Louis
(2019)}\label{e7-deslauriers-louis-2019}

Measuring actual learning versus feeling of learning in response to
being actively engaged in the classroom. \emph{Proceedings of the
National Academy of Sciences} 116(39), 19251--19257. DOI:
\url{https://doi.org/10.1073/pnas.1821936116}.

\textbf{확인한 근거.} 대학 기초 물리에서 동일 자료를 사용한 무작위
비교는 실제 평가 성과와 스스로 느낀 학습감이 다른 방향일 수 있음을
보였다. 어려운 문장이 더 좋은 수업이라는 결론은 아니다.

\textbf{이 강의에서의 선택.} 읽기 경험과 독립 수행을 따로 묻는다.
어렵다는 피드백을 무시하지 않고 모호한 문장·과도한 선수지식·설치 문제와
필요한 사고 노력을 분리한다.

\subsection{{[}E8{]} \{van Gog\}, Tamara
(2012)}\label{e8-van-gog-tamara-2012}

A Test of the Testing Effect: Acquiring Problem-Solving Skills From
Worked Examples. \emph{Cognitive Science} 36(8), 1532--1541. DOI:
\url{https://doi.org/10.1111/cogs.12002}.

\textbf{확인한 근거.} 과학 선수지식이 제한적인 대학생 40명의 회로
문제에서 예제만 공부한 조건과 예제 뒤 유사 문제를 푼 조건을 비교했다.
즉시 성과는 비슷했고 일주일 뒤에는 예제 학습 조건이 더 좋았다.

\textbf{이 강의에서의 선택.} 새롭고 복잡한 Cowling 유도에는 설명과
완전한 풀이를 충분히 남긴다. 스스로 해 보라는 요구를 모든 학습 단계의
동일한 처방으로 쓰지 않는다. 이 회로 연구도 숙련된 석사에게 같은 효과를
보장하지 않는다.

\end{document}
